\documentclass{article}
\usepackage[T1]{fontenc}
\usepackage{lmodern}
\usepackage{graphicx}
\usepackage{amsmath,amssymb}
\usepackage[a4paper,margin=2cm]{geometry}
\usepackage[absolute,overlay]{textpos}
\setlength{\TPHorizModule}{1mm}
\setlength{\TPVertModule}{1mm}
\usepackage[indonesian]{babel}
\usepackage{hyperref}
\setlength{\emergencystretch}{2em}
% unit: Halaman 18

\begin{document}

% === Halaman 18 (python/native) ===

(i) Enkripsi m1 = 0700

Bob memilih bilangan acak k = 1463 (nilai k masih berada di dalam selang [0, 2273 – 1]).

Bob mengenkripsi pesan dengan menggunakan kunci publik Alice (y = 461, g = 3, p = 2273):

a = gk mod p = 31463 mod 2273 = 1439

b = ykm1 mod p = 4611463 700 mod 2273 = 74

Jadi, cipherteks yang dihasilkan untuk m1 adalah c1 = (1439, 74).

(ii) Enkripsi m2 = 1114

Bob memilih bilangan acak k = 2001 (nilai k masih berada di dalam selang [0, 2273 – 1]). 

Bob mengenkripsi pesan dengan menggunakan kunci publik Alice (y = 461, g = 3, p = 2273):

a = gk mod p = 32001 mod 2273 = 1220

b = ykm2 mod p = 4612001 1114 mod 2273 = 1682

Jadi, cipherteks yang dihasilkan untuk m2 adalah c2 = (1220, 1682).

Bob mengirim cipherteks (1439, 74) dan (1220, 1682) kepada Alice.

18

\end{document}
